\section*{Chapitre \ref{ch:b2:pt-paths} --- Thermobarométrie et chemins P--T--t}

\begin{solution}{exo:b2:pt-paths:1}
(a) A zone à chlorite, B zone à biotite, C zone à grenat, D zone à staurolite, E zone à disthène,
F zone à sillimanite. (b) Cinq~: les isogrades de la biotite, du grenat, de la staurolite, de la disthène et de la sillimanite,
chacun entre deux échantillons successifs, le grade montant de A à F. (c) Un isograde
indique seulement qu'une réaction a eu lieu~; un gradient de terrain demande des pressions et des températures, qui
viennent de réactions calibrées ou de la thermobarométrie de l'assemblage au pic de chaque
roche.
\end{solution}

\begin{solution}{exo:b2:pt-paths:2}
(a) $(X_{\text{Mg}}/X_{\text{Fe}})_{\text{grt}} = 0.10/0.80 = 0.125$,
$(X_{\text{Mg}}/X_{\text{Fe}})_{\text{bt}} = 0.45/0.50 = 0.900$; $K_D = 0.1389$, $\ln K_D =
-1.9741$. (b) $3R = 24.942$~; dénominateur $20.19 + 24.942\times1.9741 = 69.428$~; à
\qty{5}{kbar}, $T = 54185/69.428 = \qty{780.5}{K} = \qty{507}{\celsius}$~; à
\qty{7}{kbar}, $T = (52630 + 2177)/69.428 = \qty{789.4}{K} = \qty{516}{\celsius}$.
(c) Deux kilobars la modifient de seulement \qty{9}{K}~: un thermomètre.
\end{solution}

\begin{solution}{exo:b2:pt-paths:3}
(a) $z = P/(\bar\rho g) = 6.0\times10^8/(2800\times9.81) = \qty{21.8}{km}$.
(b) $(580 - 10)/21.8 = \qty{26}{\celsius\per\km}$. (c) À \qty{580}{\celsius}, la droite
se situe à $3.67 + 0.8\times2.10 = \qty{5.35}{kbar}$~; à \qty{6.0}{kbar}, la roche est du côté
des hautes pressions~: disthène.
\end{solution}

\begin{solution}{exo:b2:pt-paths:4}
(a) Réactifs $2\times140.693 + 3\times18.015 = \qty{335.431}{g}$~; produits $277.112 +
58.320 = \qty{335.432}{g}$~: équilibré à l'arrondi près. (b) $2\times43.65 =
\qty{87.30}{cm^3}$ deviennent $107.50 + 24.63 = \qty{132.13}{cm^3}$~: $+\qty{51.4}{\percent}$.
(c) Olivine $281.386/87.30 = \qty{3.223}{g.cm^{-3}}$~; produits $335.432/132.13 =
\qty{2.539}{g.cm^{-3}}$. (d) La roche gonfle d'environ la moitié de son volume et devient plus légère~; elle
aussi stocke de l'eau dans ses minéraux.
\end{solution}

\begin{solution}{exo:b2:pt-paths:5}
(a) Moyennes $-1.206$ à \qty{1072.15}{K} et $-1.782$ à \qty{823.15}{K}. Avec $1/T =
9.3271\times10^{-4}$ et $1.21485\times10^{-3}$, $a = (-1.206 + 1.782)/(9.3271\times10^{-4}
- 1.21485\times10^{-3}) = -2041.5$ K et $b = -1.206 + 2041.5/1072.15 = 0.6982$.
(b) $\Delta_r H + 2070\times0.311 = 2041.5\times24.942 = \qty{50920}{J}$, donc $\Delta_r H =
\qty{50276}{J}$~; $\Delta_r S = 0.6982\times24.942 = \qty{17.41}{J.K^{-1}}$. (c) À
\qty{5}{kbar}~: pour $K_D = 0.18$, 588.1 contre \qty{587.5}{\celsius}~; pour $0.10$, 419
contre \qty{425}{\celsius}~; pour $0.30$, 819 contre \qty{806}{\celsius}. Dans l'intervalle
des expériences, les deux calibrations concordent, bien que leurs constantes diffèrent de
\qty{2.4}{kJ} et \qty{2.8}{J.K^{-1}}~; les erreurs de $\Delta_r H$ et $\Delta_r S$
se compensent là, et croissent en dehors. Une calibration ne doit pas être utilisée loin
des conditions de ses expériences.
\end{solution}

\begin{solution}{exo:b2:pt-paths:6}
(a) Échange~: $(113.12 + 154.30) - (115.32 + 149.65) = \qty{+2.45}{cm^3} =
\qty{0.245}{J.bar^{-1}}$, contre 0.311. GASP~: $125.28 + 2\times44.15 + 22.69 -
3\times100.79 = \qty{-66.10}{cm^3} = \qty{-6.61}{J.bar^{-1}}$, contre $-6.21$, la valeur
de la calibration à \qty{800}{\celsius} et \qty{10}{kbar}, où la dilatation thermique et
la compression modifient les volumes. Dans les deux cas, la variation GASP est vingt à trente fois
celle de l'échange. (b) $\partial T/\partial P = 0.311/62.96 = \qty{4.9e-3}{K/bar}$~:
\qty{4.9}{K} par kilobar. (c) La pente de la droite vaut $22.80 + 8.314\times3\ln 0.20/6.21
= 22.80 - 6.46 = \qty{16.3}{bar/K}$~; \qty{17}{K} décalent la pression de
\qty{280}{bar}.
\end{solution}

\begin{solution}{exo:b2:pt-paths:7}
(a) $P_0 = 22.80\times600 - 1093 = \qty{12587}{bar}$~; $3\ln(0.08/0.30) = -3.9653$~; $RT =
8.314\times873.15 = \qty{7259.4}{J}$; $P = 12587 + 7259.4\times(-3.9653)/6.21 = 12587 -
4635 = \qty{7952}{bar}$. (b) $3\ln(0.08/0.35) = -4.4277$, $P = 12587 - 5176 =
\qty{7411}{bar}$~: \qty{0.54}{kbar} de moins, donc le plagioclase doit être analysé avec soin.
(c) $7952\times10^5/(2800\times9.81) = \qty{29.0}{km}$.
\end{solution}

\begin{solution}{exo:b2:pt-paths:8}
(a) $z = P/(\bar\rho g)$~: \qty{118.9}{km}, \qty{34.0}{km} et \qty{15.3}{km}.
(b) $(118.9 - 34.0)/2.2 = \qty{38.6}{km/Ma} = \qty{3.9}{cm.an^{-1}}$ et $(34.0 -
15.3)/1.1 = \qty{17.0}{km/Ma} = \qty{1.7}{cm.an^{-1}}$, un peu au-dessus des valeurs publiées
3.4 et 1.6~: leur conversion de pression en profondeur donne des écarts de profondeur un peu plus petits
différences. (c) $\sigma = \sqrt{0.9^2 + 0.9^2} = \qty{1.3}{Ma}$, donc $2.2 \pm 1.3$
millions d'années et une vitesse comprise entre $84.9/3.47 = 24$ et $84.9/0.93 = \qty{91}{km/Ma}$,
2.4 à \qty{9.1}{cm.an^{-1}}~: rapide dans tous les cas. (d) Même la vitesse la plus faible est environ cinquante fois
les \qty{0.5}{mm.an^{-1}} du modèle~: la seule érosion ne peut pas y suffire.
\end{solution}

\begin{solution}{exo:b2:pt-paths:9}
(a) $\partial T/\partial\ln K_D = 3RT/(\Delta_r S - 3R\ln K_D) = 24.942\times860.6/62.96 =
\qty{341}{K}$~; $\partial T/\partial P = 0.311/62.96 = \qty{4.9}{K}$ par kilobar.
(b) $\sigma_T = \sqrt{(341\times0.05)^2 + 4.9^2} = \sqrt{17.0^2 + 4.9^2} =
\qty{17.7}{K}$. (c) L'erreur analytique sur $K_D$ domine~; à elle seule elle représente les deux tiers de
les \qty{25}{\celsius} des meilleures calibrations, qui incluent aussi les erreurs de la
calibration elle-même.
\end{solution}

\begin{solution}{exo:b2:pt-paths:10}
(a) Grenat $0.10/0.75 = 0.1333$, biotite $0.42/0.52 = 0.8077$, $K_D = 0.1650$, $\ln K_D =
-1.8018$, dénominateur $20.19 + 24.942\times1.8018 = 65.130$~; $3\ln(0.07/0.25) = -3.8189$.
À \qty{6000}{bar}~: $T = 54496/65.130 = \qty{836.7}{K} = \qty{563.6}{\celsius}$, puis $P
= 22.80\times563.6 - 1093 + 8.314\times836.7\times(-3.8189)/6.21 = \qty{7479}{bar}$.
À \qty{7479}{bar}~: $T = \qty{570.6}{\celsius}$, $P = \qty{7603}{bar}$. À
\qty{7603}{bar}~: $T = \qty{571.2}{\celsius}$, $P = \qty{7614}{bar}$~; une fois de plus, $T =
\qty{571.3}{\celsius}$. La roche s'est équilibrée vers \qty{571}{\celsius} et
\qty{7.6}{kbar}. (b) À \qty{571}{\celsius}, la droite disthène $=$ sillimanite se situe à $3.67
+ 0.71\times2.10 = \qty{5.16}{kbar}$~: disthène. (c) Le baromètre a été calibré pour
la disthène~; avec la sillimanite, la réaction et sa droite diffèrent, car les deux polymorphes n'ont pas
la même énergie de Gibbs.
\end{solution}

\begin{solution}{exo:b2:pt-paths:11}
(a) $\tau = L^2/(\pi^2\kappa)$~: pour \qty{35}{km}, $1.225\times10^9/(9.870\times10^{-6}) =
\qty{1.24e14}{s} = \qty{3.9}{Ma}$~; pour \qty{70}{km}, quatre fois plus,
\qty{15.7}{Ma}. (b) Soixante-dix millions d'années représentent quatre à dix-huit fois ces temps~: la chaleur
traverse la croûte épaissie bien plus vite que l'érosion ne l'enlève. (c) La roche
atteint sa plus grande profondeur d'un coup, lors de l'empilement~; la chaleur arrive pendant les
dizaines de millions d'années suivantes, tandis que l'érosion la remonte lentement~: la température continue de
monter un temps pendant la décompression. (d) Une érosion en sept millions d'années ferait
remonter la roche avant que la conduction ne l'ait réchauffée~: un pic thermique plus bas, atteint
tôt, à plus haute pression, et un chemin froid jusqu'à la surface.
\end{solution}

\begin{solution}{exo:b2:pt-paths:12}
(a) La croûte pèse $35000\times2800\times9.81 = \qty{0.961}{GPa}$~; les
\qty{1.839}{GPa} restants demandent $1.839\times10^9/(3300\times9.81) = \qty{56.8}{km}$ de manteau~: profondeur
\qty{91.8}{km}, comme les \qty{90}{km} ou plus de l'étude originale. (b) $(4.0 -
0.961)\times10^9/(3300\times9.81) = \qty{93.9}{km}$, profondeur \qty{128.9}{km}.
(c) $700/91.8 = \qty{7.6}{\celsius\per\km}$~: un rapport pression sur température très élevé, la série
des zones de subduction. (d) $(91.8 - 35)/34 = \qty{1.7}{Ma}$.
\end{solution}

\begin{solution}{pb:b2:pt-paths:1}
\noindent\textbf{Partie I.}

\textbf{1.} Cœur de grenat $0.080/0.760 = 0.1053$, biotite $0.40/0.58 = 0.6897$~; $K_D =
0.1527$, $\ln K_D = -1.8793$.

\textbf{2.} Bordure $0.120/0.720 = 0.1667$, biotite de la matrice $0.46/0.53 = 0.8679$~; $K_D =
0.1921$, $\ln K_D = -1.6497$.

\textbf{3.} $3R = 24.942$~; numérateur $52630 + 0.311\times6000 = 54496$~; dénominateur $20.19
+ 24.942\times1.8793 = 67.064$~; $T = \qty{812.6}{K} = \qty{539.5}{\celsius}$.

\textbf{4.} Dénominateur $20.19 + 24.942\times1.6497 = 61.337$~; $T = 54496/61.337 =
\qty{888.5}{K} = \qty{615.3}{\celsius}$.

\textbf{5.} $0.311\times1000/67.064 = \qty{4.6}{K}$ pour le cœur et $0.311\times1000/61.337
= \qty{5.1}{K}$ pour la bordure.

\textbf{6.} Un thermomètre s'applique à des minéraux qui étaient en équilibre ensemble. La
biotite incluse dans le cœur a été isolée du reste de la roche lorsque le cœur a cristallisé~;
la biotite de la matrice a échangé pour la dernière fois avec la bordure. Associer le cœur à la matrice mélangerait
deux moments.

\noindent\textbf{Partie II.}

\textbf{7.} Cœur $22.80\times539.5 - 1093 = \qty{11208}{bar}$~; bordure $22.80\times615.3 -
1093 = \qty{12936}{bar}$.

\textbf{8.} $3\ln(0.050/0.25) = -4.8283$~; $RT = 8.314\times812.65 = \qty{6756.4}{J}$~; $P =
11208 - 6756.4\times4.8283/6.21 = 11208 - 5253 = \qty{5955}{bar}$.

\textbf{9.} $3\ln(0.060/0.28) = -4.6213$~; $RT = 8.314\times888.45 = \qty{7386.6}{J}$~; $P
= 12936 - 5497 = \qty{7439}{bar}$.

\textbf{10.} À \qty{7439}{bar}~: $T = (52630 + 2314)/61.337 = \qty{895.8}{K} =
\qty{622.6}{\celsius}$~; à cette température $P = 13102 - 5542 = \qty{7560}{bar}$~; puis $T
= \qty{623.2}{\celsius}$. La bordure a cristallisé vers \qty{623}{\celsius} et \qty{7.56}{kbar}.

\textbf{11.} À \qty{5954}{bar}, $T = \qty{539.2}{\celsius}$ et $P = \qty{5950}{bar}$~: le
cœur a cristallisé vers \qty{539}{\celsius} et \qty{5.95}{kbar}. La première supposition,
\qty{6}{kbar}, était déjà proche, et le thermomètre ne bouge que de \qty{4.6}{K} par kilobar.

\textbf{12.} Cœur $5950\times10^5/(2800\times9.81) = \qty{21.7}{km}$~; bordure
$7560\times10^5/27468 = \qty{27.5}{km}$.

\textbf{13.} À \qty{539}{\celsius}, la droite se situe à $3.67 + 0.39\times2.10 =
\qty{4.49}{kbar}$, sous \qty{5.95}{kbar}~: disthène, compatible avec la disthène incluse
dans le cœur. À \qty{623}{\celsius}, elle se situe à $5.77 + 0.23\times2.09 = \qty{6.25}{kbar}$,
sous \qty{7.56}{kbar}~: de nouveau la disthène.

\noindent\textbf{Partie III.}

\textbf{14.} $+\qty{84}{\celsius}$ et $+\qty{1.6}{kbar}$ (environ \qty{6}{km} plus profond)~: la
roche était enfouie et réchauffée, sur son \omterm{def:b2:pt-paths:paths}{chemin prograde}.

\textbf{15.} À \qty{623}{\celsius} ou plus, sous \qty{6.25}{kbar}~: après la croissance de la bordure, la
roche a perdu au moins \qty{1.3}{kbar}, environ \qty{5}{km} de couverture, tout en restant au moins
aussi chaude. Elle a été exhumée encore chaude.

\textbf{16.} La chlorite est un minéral hydraté de bas grade~: elle enregistre le \omterm{def:b2:pt-paths:paths}{chemin rétrograde} tardif,
qui se refroidit à faible profondeur, en présence d'eau.

\textbf{17.} Enfouissement et réchauffement jusqu'à la bordure, décompression sans refroidissement, puis refroidissement~:
tracé avec la température vers la droite et la pression vers le haut, le chemin tourne dans le sens horaire. Le
pic de pression est survenu le premier, à la bordure~; le pic thermique au même moment ou plus tard, pendant la
décompression.

\textbf{18.} Oui~: le modèle atteint aussi son pic de pression en premier et son pic thermique,
\qty{593}{\celsius} à \qty{7.8}{kbar}, pendant la remontée, proches des conditions de la bordure
du schiste. L'épaississement suivi d'une érosion lente explique le chemin du schiste.

\noindent\textbf{Partie IV.}

\textbf{19.} $(623 - 405)/(412.0 - 398.0) = 218/14.0 = \qty{15.6}{\celsius}$ par million
d'années.

\textbf{20.} La roche s'est refroidie environ une fois et demie plus vite que supposé. Un minéral
qui se refroidit plus vite a moins de temps pour perdre son argon, donc il se ferme à une température plus élevée~:
la vraie température de fermeture se situe un peu au-dessus de \qty{405}{\celsius}, et la fermeture
a lieu un peu plus profondément.

\textbf{21.} $(405 - 10)/30 = \qty{13.2}{km}$.

\textbf{22.} $27.5 - 13.2 = \qty{14.3}{km}$ (avec les profondeurs non arrondies, 27.52 et 13.17,
\qty{14.35}{km}).

\textbf{23.} $\sqrt{1.5^2 + 1.0^2} = \qty{1.8}{Ma}$~: $14.0 \pm \qty{1.8}{Ma}$, donc une vitesse
comprise entre $14.35/15.8 = 0.91$ et $14.35/12.2 = \qty{1.18}{km/Ma}$.

\textbf{24.} $14.35/14.0 = \qty{1.03}{km}$ par million d'années, soit
\textbf{\qty[text-series-to-math=true]{1.0}{mm.an^{-1}}}, environ trente fois plus lent que
le premier stade d'\omterm{def:b2:pt-paths:exhumation}{exhumation} des roches alpines les plus profondes, et environ le double de la vitesse d'érosion
du modèle d'épaississement.
\end{solution}
